\documentclass[11pt]{article}
\usepackage{amsmath,amssymb}
\usepackage[hidelinks]{hyperref}
\title{A restart that preserves an accepted right, and the cost of announcing it}
\author{}\date{}
\begin{document}\maketitle

\paragraph{Lift and scope.}
The protected obligation is an old accepted return through a shared bay,
not just a previous step's receipt nor a future robot position. The
new visitor needs a lawful entry, and other work on the same timeline
must continue. The source roles and their permissions are authored for
a civil development fixture. Native RWARE supplies the selected motion;
native SubT supplies actual packet charges and deliveries. An authored
SQLite authority queue persists acceptance independently of executing the
motion. This does not attribute such a queue or a physical interlock to
stock RWARE or to hardware.

\paragraph{Paired restart.}
A subprocess exits around its authority commit. When it exits before
commit, the return is absent after rollback; when it exits after commit,
the named return is accepted but unexecuted. The two public restart
views are identical: reset, old carrier at \((3,0)\), shared bay empty.
A local transaction seals epoch 7 before the source reads the named
right. This forbids new issue in the fixture, not retained execution.
No surviving writer, unrecorded alternate issuer or pre-existing
unaudited command is in this envelope. The checked read concerns only
the same epoch and key, not general absence of motion.

A qualified negative certificate can admit the visitor at tick 6.
If the old right is present, its actual carrier must first move into
and then out of \((2,0)\). Only the actual motion and subsequent
clearance allow the typed downstream release. Merely granting the
visitor access because the bay looks empty instead allows its native
move first, after which native occupancy blocks the protected carrier.
In this diagnostic arm the conflict fails the old service, not a
physical collision. Tick 5 also carries an actual disjoint transmission
and a separate native carrier move while the affected route is down.

\paragraph{Create the evidence that changes a decision.}
The first setup reaches \(Q\) at tick 4. The next fault can erase its
proof transmission at tick 5 after setup. A preauthorized alternative
\(X\) can carry the same qualified observation at tick 6 if available.
Silence from either route never means that the old job is absent.
The source sends the early 3-byte certificate only when it can enable
early entry. Announcing the presence of a still-live return is truthful,
but in this fixture it does not advance a positive action. If that
announcement is received at tick 5 or 6, the actual 3-byte release
at tick 9 needs capacity 6 at the shared receiver. At capacity 5 it
fails; the send attempt charges the release sender anyway. At tick 11
the early announcement is gone from the receiver window. The test
uses \(X\) at that point to avoid assuming that the tick-9 sender debit
has already expired. A late link, no backup, or no early receipt can
still produce lawful entry through the reserved release at tick 9.

\paragraph{Comparators.}
The guard preserves both obligations by tick 9 in all 16 enumerated
histories and allows 6 qualified early entries. A safe static wait
preserves all 16 but allows no early entry. Eager truthfully publishing
both kinds of early report matches the six early entries, but leaves
three later cases beyond tick 9. The occupancy-only diagnostic loses
the eight retained old duties. The equally informed direct calculation
matches all 64 policy histories. The difference is a change to
evidence transmission and its capacity use, not a stronger optimization
algorithm or a measured failure rate. The early fallback availability,
freshness of the named read, stable seal and late release link all
remain explicit premises. If a live-right announcement were itself
needed for safety, fairness or service elsewhere, the suppression rule
would not transport without revising that obligation. No full-system
runtime, independent mission or I1 acceptance.
\end{document}
